% Part: first-order-logic
% Chapter: syntax
% Section: intro-syntax

\documentclass[../../../include/open-logic-section]{subfiles}

\begin{document}

\olfileid{fol}{syn}{itx}

\olsection{引言}

为了发展一阶逻辑的理论与元理论，我们必须首先定义其表达式的语法与语义。
一阶逻辑的表达式是项与公式。
项由变元、常项符号与函数符号构成。
而公式则由谓词符号与项构成（它们构成最小的“原子”公式），进而我们可以利用逻辑联结词与量词，从原子公式构成更复杂的公式。
规定形成规则的方式有很多种；我们只给出其中一种。
其他系统会选择不同的符号，选取不同的联结词集合作为初始联结词，并以不同方式使用括号（甚至完全不用括号，如所谓的波兰记法）。
不过，所有方法的共同之处在于，形成规则都\emph{归纳地}定义了项与公式的集合。
如果定义得当，每个表达式本质上只能由形成规则以一种方式生成。
由于这种归纳定义生成的表达式是\emph{唯一可读的}，这意味着我们可以使用相同的方法——归纳定义——来为这些表达式赋予意义。

\end{document}